%% ============================================================
%% 分片 12：原文第 5 节「Topological Field Theory In Two Dimensions」（原文 glshort.tex 起于第 3264 行）
%% 由 tools/harvmac2tex.py 从 arXiv 源码（hep-th/0604151v3）自动生成。
%% 只译英文散文；公式、编号、\label、\cite 一律不改。
%% ============================================================

\subsection{Twisted Topological Field Theories}

\label{sec:twistedtft}

We now want to study the sigma-model of target $\MH$ from the
viewpoint of topological field theory, so we replace
$\Bbb{R}^{1,1}$ by a general two-manifold $\Sigma$ that we
eventually will take to have Euclidean signature. We want to
discuss what twisted topological field theories on $\Sigma$ can be
constructed from the sigma-model of target $\MH$, and which of
these actually arise by compactifying on the Riemann surface $C$
the four-dimensional topological field theories that we
constructed in section 第~\ref{sec:tqftfour}~节.

 
 

 

 
 

To construct a topological field theory from a sigma-model with a
target $X$, one picks a {\it pair} of complex structures
$(J_+,J_-)$ on $X$ that obey certain conditions of $(2,2)$
supersymmetry \cite{ref66}.\footnote{Generalized complex geometry leads to an
extension of this construction, which we will not need here, in
which only one of the two generalized complex structures that can
be formed \cite{ref67}\ from the pair $(J_+,J_-)$ is integrable.
One then aims as in \cite{ref68}\ to define a topological field theory using
this  integrable generalized complex structure.}
These conditions have recently
been reinterpreted \cite{ref67}\ in terms of generalized complex
geometry \cite{ref69}. Once one has $(2,2)$ supersymmetry, there is a
standard recipe \cite{ref70}, via the twisting procedure that we
reviewed in section 第~\ref{sec:twisting}~节, for constructing a topological
field theory.

If $X$ is hyper-Kahler, things simplify. The sigma-model with
target $X$ has $(4,4)$ supersymmetry, and a structure of $(2,2)$
supersymmetry can be obtained by picking a $(2,2)$ subalgebra of
the $(4,4)$ supersymmetry algebra.\footnote{Two $(2,2)$ subalgebras
that differ by conjugation by $Spin(4)_{\cal R}$ are considered
equivalent.} In terms of choosing complex structures, this amounts
to the following.  $X$ has a family of complex structures
parametrized by a copy of $\Bbb{CP}^1$ that we call
$\Bbb{CP}^1_h$, and a $(2,2)$ structure can be defined by picking
a pair of points $(J_+,J_-)\in \Bbb{CP}_h^1$. The conditions in
\cite{ref66}\ are automatically obeyed for such a pair. In the case of
$\MH$, this gives the family of topological field theories that is
important for the geometric Langlands program.

We can conveniently characterize as follows the topological field
theory associated with the pair $(J_+,J_-)$.  Let $\Phi: \Sigma\to
\MH$ be the bosonic field of the sigma-model.   Then the equations
of unbroken supersymmetry read
% 式 (5.5)
\setcounter{equation}{4}
\begin{equation}\label{eq:hson}
(1-iJ_+)\bar\partial\Phi=0,
\end{equation}
 and
% 式 (5.6)
\setcounter{equation}{5}
\begin{equation}\label{eq:hsono}
(1-iJ_-)\partial \Phi=0.
\end{equation}
  Here $\partial$ and
$\bar\partial$ are the usual $(1,0)$ and $(0,1)$ parts of the
exterior derivative on $\Sigma$; and $\partial\Phi$ and
$\bar\partial \Phi$ are understood as one-forms on $\Sigma$ with
values in the pullback of the tangent bundle $TX$ (on which $J_+$
and $J_-$ act). The first condition says that the map
$\Phi:\Sigma\to X$ is holomorphic in complex structure $J_+$; the
second says that it is antiholomorphic in complex structure $J_-$.
Analogously to how we found the conditions of unbroken
supersymmetry in section 第~\ref{sec:familya}~节, these equations arise because
there are fermi fields $\chi_+$, $\chi_-$ with
$\{Q,\chi_+\}=(1-iJ_+)\bar\partial\Phi$,
$\{Q,\chi_-\}=(1-iJ_-)\partial\Phi$.

There is much more information in the equations \eqref{eq:hson}\ and \eqref{eq:hsono}\
than in their solutions. To discuss the solutions of the
equations, one usually specializes to $\Sigma$ of Euclidean
signature (the most natural case for topological field theory).
Then for generic $J_+$, $J_-$,  they imply that the map
$\Phi:\Sigma\to X$ must be constant, a result that does not depend
on $J_+$ and $J_-$.  The equations themselves, understood
algebraically, and without fixing a real structure on the
two-manifold $\Sigma$, depend on $J_+$ and $J_-$.

An important example is the case $J_+=J_-=\hat J$ (where, in our
applications, $\hat J$ will be one of the complex structures $I_w$
of $\MH$, as described in eqn. \eqref{eq:wyre}). The equations combine to
$(1-i\hat J\,)d\Phi = 0$, which imply that $d\Phi=0$. The model is
called the $B$-model in complex structure $\hat J$. If instead
$J_+=-J_-=\hat J$, we get the $A$-model in complex structure $\hat
J$. (This model really depends on the symplectic form
$\omega_{\hat J}$ associated with complex structure $\hat J$
rather than  the complex structure {\it per se}. But calling it
the $A$-model in complex structure $\hat J$ is sometimes a
convenient shorthand.) For $\Sigma$ of Euclidean signature, the
equations of the $A$-model are redundant, as eqn. \eqref{eq:hsono}\ is the
complex conjugate of \eqref{eq:hson}. Either one of them asserts that the
map $\Phi:\Sigma\to \MH$ is holomorphic in complex structure $\hat
J$. If $J_-\not= \pm J_+$, we get a model that is neither an
$A$-model nor a $B$-model, but can be reduced to an $A$-model
using generalized complex geometry, as we discuss in section
第~\ref{sec:gencpx}~节.

We recall from eqn. \eqref{eq:wyre}\ the explicit form of the family of
complex structures on $\MH$ parametrized by $\Bbb{CP}_h^1$.  In
terms of an affine parameter $w$ for $\Bbb{CP}_h^1$, we defined a
complex structure $I_w$ in which the holomorphic coordinates are
$A_{\bar z}-w\phi_{\bar z}$ and $A_z+w^{-1}\phi_z$. Explicitly, we
found % 式 (5.7)
\setcounter{equation}{6}
\begin{equation}\label{eq:wyrre}
I_w={1-\bar w w\over 1+\bar w w}I+{i(w-\bar
w)\over 1+\bar w w}J+{w+\bar w\over 1+\bar w w}K.
\end{equation}

 The pair
$(J_+,J_-)$ corresponds in this parametrization to a pair
$(w_+,w_-)$.

\par\medskip{\it Reduction From Four Dimensions}

Now let us compare this to what we get by reduction from four
dimensions.  We learned in section 第~\ref{sec:tqftfour}~节\ that gauge theory
leads to a family of topological quantum field theories (TQFT's)
defined on any four-manifold $M$ and parametrized by a copy of
$\Bbb{CP}^1$ that we will call $\Bbb{CP}^1_g$.   This family
certainly reduces for $M=\Sigma\times C$ to a subfamily of the
family $\Bbb{CP}_h^1\times \Bbb{CP}_h^1$ that is natural from the
two-dimensional point of view. Both families come from a twisting
procedure applied to ${\EUN}=4$ super Yang-Mills theory on a
four-manifold $M$. In this procedure, a topological field theory
is obtained from the cohomology of a suitable linear combination
$Q$ of the supersymmetries.  To get the family $\Bbb{CP}^1_g$, we
require $Q$ to be $Spin'(4)$-invariant so that the construction
works for all $M$. But to get the sigma-model of target $\MH$, we
specialize to the manifold $M=\Sigma\times C$ of reduced holonomy,
and then, as we discuss more explicitly later, a weaker condition
on $Q$ suffices, leading to the larger family $\Bbb{CP}_h^1\times
\Bbb{CP}_h^1$.

We aim to identify the four-dimensional topological field theory
family $\Bbb{CP}^1_g$ as a subfamily of the sigma-model family
$\Bbb{CP}_h^1\times \Bbb{CP}_h^1$. The embedding of $\Bbb{CP}^1_g$
in $\Bbb{CP}_h^1\times \Bbb{CP}_h^1$ can be described by functions
$w_+(t)$, $w_-(t)$. We can compute these functions by specializing
to a convenient configuration or physical state.

A very convenient way to proceed is to simply abelianize the
problem, working in a vacuum in which $G$ is broken to its maximal
torus by the expectation values of some of the untwisted scalar
fields.   In the abelian case, from \eqref{eq:loopme}, the equations of
unbroken supersymmetry in four dimensions are
% 式 (5.8)
\setcounter{equation}{7}
\begin{equation}\label{eq:onso}
\begin{aligned}(F+t\,d\phi)^+&=0\\(F-t^{-1}d\phi)^-&=0.\\\end{aligned}
\end{equation}

Here, the one-form $\phi$ is a section of $T^*M=T^*\Sigma\oplus
T^*C$, and the gauge field $A$ is locally a one-form or section of
$T^*M$. To compare to a sigma-model on $\Sigma$ with target
$\MH(G,C)$, we must take $\Sigma $ to have radius much greater
than that of $C$.  In that limit, the parts of $A$ and $\phi$ that
take values in $T^*\Sigma$ are ``massive'' and can be dropped; so
$A$ and $\phi$ can be interpreted as sections of $T^*C$, slowly
varying on $\Sigma$.

We write $y$ for a local complex coordinate on $\Sigma$ and $z$
for a local complex coordinate on $C$.  The two-form $d\bar
y\wedge d\bar z$ is selfdual, while $dy\wedge d\bar z$ is
anti-selfdual. So the first equation in \eqref{eq:onso}\ leads to
% 式 (5.9)
\setcounter{equation}{8}
\begin{equation}\label{eq:igot}
\partial_{\bar y}\left(A_{\bar z}+t \phi_{\bar
z}\right)=0,
\end{equation}
 and the second leads to
% 式 (5.10)
\setcounter{equation}{9}
\begin{equation}\label{eq:nigot}
\partial_y\left(A_{\bar z}-t^{-1}\phi_{\bar
z}\right)=0.
\end{equation}

By comparing these results to eqns. \eqref{eq:hson}\ and \eqref{eq:hsono}, we can read
off the functions $w_+(t)$ and $w_-(t)$.   Eqn. \eqref{eq:igot}\ asserts
that the map $\Sigma\to \MH(G,C)$ is holomorphic if we take on
$\MH(G,C)$  the complex structure in which $A_{\bar z}+t\phi_{\bar
z}$ is holomorphic.  This agrees with eqn. \eqref{eq:hson}, which asserts
such holomorphy in complex structure $I_{w_+}$, if and only if
$w_+=-t$. Similarly, eqn. \eqref{eq:nigot}\ asserts that the map $\Sigma\to
\MH(G,C)$ is antiholomorphic if we take on $\MH(G,C)$ the complex
structure in which $A_{\bar z}-t^{-1} \phi_{\bar z}$ is
holomorphic.  This agrees with eqn. \eqref{eq:hsono}, which asserts
antiholomorphy in complex structure $I_{w_-}$, if and only if
$w_-=t^{-1}$.

So the embedding of $\Bbb{CP}^1_g$ in $\Bbb{CP}_h^1\times
\Bbb{CP}_h^1$ is defined by % 式 (5.11)
\setcounter{equation}{10}
\begin{equation}\label{eq:tyroy}
\begin{aligned}w_+ & = -t\\
                                              w_- & = t^{-1}.\\\end{aligned}
\end{equation}

\par\medskip{\it Some Simple Considerations}

Let us now make a few simple simple observations about this
result.

First of all, when do we get a $B$-model?  For a $B$-model, we
want $w_+=w_-$.  So the condition is $t^{-1}=-t$, which occurs
precisely for $t=\pm i$.  Since the complex structures $I_w$ for
$w=\mp i$ coincide with $J$ and $-J$, we get this way the
$B$-model in complex structure $J$ or $-J$.

When do we get an $A$-model?  For an $A$-model, the complex
structures $I_{w_+}$ and $I_{w_-}$ should be opposite.  The map
$w\to -1/\bar w$ maps $I_w$ to its opposite (this is clear from
\eqref{eq:wyrre}), so we get an $A$-model in two dimensions if $w_-=-1/\bar
w_+$, which works out to $t=\bar t$.  In other words, precisely
for real $t$, we get an $A$-model.  For example, for $t=1$ or
$-1$, we get the $A$-model in complex structure $K$ or $-K$.  For
$t=0$ or $\infty$, we get the $ A$-model in complex structure $I$
or $-I$.

Now, let us compare this family to the more complete family
$\Bbb{CP}_h^1\times \Bbb{CP}_h^1$ of the sigma-model.  At the end
of section 第~\ref{sec:complexstr}~节, we described the group ${\cal U}\cong
\Bbb{C}^*$ of diffeomorphisms of $\MH(G,C)$; it acts on the family
of complex structures $I_w$ by $w\to \lambda^{-1} w$, $\lambda\in
\Bbb{C}^*$.   The topological field theory determined by a
pair\footnote{This topological field theory depends also on the
hyper-Kahler metric and $B$-field of $\MH$.  The full dependence
on all variables is determined in section 第~\ref{sec:gencpx}~节\ using
generalized complex geometry.} $(w_+,w_-)$ is the same as that
determined by $(\lambda^{-1} w_+,\lambda^{-1} w_-)$. How many
really inequivalent TQFT's can we define from the two-dimensional
point of view?  The only invariant we can form from the pair
$(w_+,w_-)$ is the ratio $q=w_+/w_-$. We see that
% [WARN] 公式 \eqn\eqref{eq:yto} 无编号或重复
\begin{equation*}
\tag*{(5.12)}
q=-t^2,
\end{equation*}
 so all values of $q$ can be achieved, but not
quite uniquely.  The points $t$ and $-t $ on $\Bbb{CP}^1_g$ lead
to equivalent theories in two dimensions. This equivalence
reflects the action of the center of $SU(4)_{\cal R}$ (recall eqn.
\eqref{eq:ungu}).

Although all values of the invariant $q$ do come from
four-dimensional TQFT's, it is not quite true that all $\Bbb{C}^*$
orbits on $\Bbb{CP}_h^1\times \Bbb{CP}_h^1$ have representatives
with such an origin.  The missing orbits are the
$\Bbb{C}^*$-invariant points  $(0,0)$ and $(\infty,\infty)$, and
also the orbits in which $w_+$ or $w_-$, but not both, is 0 or
$\infty$. Particularly notable is the fact that the points $(0,0)$
and $(\infty,\infty)$ are not equivalent to theories that
originate in four-dimensional TQFT's. These points correspond to
the $B$-models in complex structures $I$ and $-I$.

The $B$-model in complex structure $I$ has been the starting point
in mathematical efforts -- briefly surveyed in the introduction --
to interpret the geometric Langlands program in terms of the
geometry of $\MH$.  Because the Hitchin fibration is holomorphic
in complex structure $I$, the $T$-duality on the fibers of the
Hitchin fibration (whose relation to four-dimensional $S$-duality
we review soon) maps the $B$-model of complex structure $I$ to
itself, acting  on $D$-branes via the Fourier-Mukai transform that
is the starting point in the mathematical description.  Although
the point $(0,0)$ corresponding to this $B$-model is not in the
family $\Bbb{CP}_g^1$ that comes from four-dimensional TQFT's, it
is interesting that it can be infinitesimally perturbed to give
points on the $\Bbb{C}^*$ orbits corresponding to almost any point
on $\Bbb{CP}^1_g$.  We simply perturb $(0,0)$ to $(\alpha,\beta)$
for infinitesimal $\alpha$, $\beta$; the invariant $q=w_+/w_-$ is
then $q=\alpha/\beta$, and can take any value for arbitrarily
small $\alpha,\beta$.  Perhaps this fact will lead eventually to
an understanding of the geometric Langlands program based on
perturbing from the $B$-model in complex structure $I$.  Our
approach, however, will rely on the family $\Bbb{CP}^1_g$ that
comes directly from four dimensions.

Although certain $\Bbb{C}^*$ orbits, such as the point $(0,0)$, do
not arise by specializing a four-dimensional TQFT to
$M=\Sigma\times C$, this does not meant that they cannot be
described in four-dimensional gauge theory.  In section 第~\ref{sec:twisting}~节,
we obtained the family $\Bbb{CP}^1_g$ using the supersymmetry
generator % 式 (5.13)
\setcounter{equation}{12}
\begin{equation}\label{eq:klop}
\epsilon=u\epsilon_\ell+v\epsilon_\rr
\end{equation}
 that
was constrained to be $Spin'(4)$-invariant.  If we specialize to
$M=\Sigma\times C$, there is no need to ask for $Spin'(4)$
invariance.  The holonomy group of $\Sigma\times C$ is
$SO(2)\times SO(2)\subset SO(4)$, and $Spin'(2)\times Spin'(2)$
invariance is enough.  This means that we can generalize \eqref{eq:klop}\ to
% 式 (5.14)
\setcounter{equation}{13}
\begin{equation}\label{eq:blop}
\hat\epsilon=\left(u+\tilde u\,
\Gamma^*\right)\epsilon_\ell
                 +\left(v+\tilde v\,\Gamma^*\right)\epsilon_\rr,
\end{equation}

                 where
% 式 (5.15)
\setcounter{equation}{14}
\begin{equation}\label{eq:ujbyty}
\Gamma^*=\begin{cases}\Gamma_{01}&\text{in Lorentz
signature}\\ i\Gamma_{01}&\text{in Euclidean signature}\\\end{cases}
\end{equation}
 is
the operator that distinguishes the two chiralities of
two-dimensional spinors.  We will adopt Euclidean signature here,
as this is more natural for topological field theory.

The supersymmetry generators in \eqref{eq:blop}\ are precisely the ones that
obey the following conditions:
% 式 (5.16)
\setcounter{equation}{15}
\begin{equation}\label{eq:zonop}
\begin{aligned}\Gamma_{2367}\epsilon& = \epsilon \\
\Gamma_{0145}\epsilon & = \epsilon.\\\end{aligned}
\end{equation}
 The first is eqn. \eqref{eq:ytytag},
which says that the supersymmetry generated by $\epsilon$ is
unbroken by the curvature of $C$; the second says that it is
similarly unbroken by the curvature of $\Sigma$. With this more
general starting point, it is possible to get the whole family
$\Bbb{CP}_h^1\times \Bbb{CP}_h^1$ of two-dimensional TQFT's from
four-dimensional gauge theory, though not from a four-dimensional
TQFT.

It actually is convenient to rewrite \eqref{eq:blop}\ in an eigenbasis of
$\Gamma^*$:
% 式 (5.17)
\setcounter{equation}{16}
\begin{equation}\label{eq:bilop}
\hat\epsilon=\half\left(1-\Gamma^*\right)(u'\epsilon_\ell+v'\epsilon_r)
+\half\left(1+\Gamma^*\right)(u''\epsilon_\ell+v''\epsilon_r).
\end{equation}

Here $(u',v')$ and $(u'',v'')$ are, respectively, homogeneous
coordinates for the two factors of $\Bbb{CP}^1_h\times
\Bbb{CP}^1_h$.  Let $\hat\delta_T$ be the extended topological
symmetry generated by $\hat \epsilon$.  To determine the two
complex structures $(J_+,J_-)$, or equivalently, to determine the
pair $(w_+,w_-)$, we just need to compute
$\hat\delta_T\chi^+_{\bar y\,\bar z}$ and
$\hat\delta_T\chi^-_{y\bar z}$. Setting these to zero will give
the generalization of \eqref{eq:onso}. To determine the generalization of
the usual formulas from eqn. \eqref{eq:helpme}, all we need to know is that
if $\tilde\chi^+=\Gamma^{\bar y\,\bar z}\chi^+_{\bar y\,\bar z}$
and $\tilde\chi^-=\Gamma^{y\bar z}\chi^-_{y\bar z}$, then
$\tilde\chi^+\Gamma^*=\tilde\chi^+$  and
$\tilde\chi^-\Gamma^*=-\tilde\chi^-$. Using this, we get
% 式 (5.18)
\setcounter{equation}{17}
\begin{equation}\label{eq:qengo}
\begin{aligned}\hat\delta_T\chi^+_{\bar y\,\bar
z}&=u'(F-\phi\wedge\phi)_{\bar y\,\bar z}+v'(D\phi)_{\bar y\,\bar
z}\\ \hat\delta_T\chi^-_{ y \bar z}&=v''(F-\phi\wedge\phi)_{y
z}-u''(D\phi)_{yz}.\\\end{aligned}
\end{equation}
 If, therefore, $t'=v'/u'$ and
$t''=v''/u''$, then the generalizations of eqns. \eqref{eq:igot}\ and
\eqref{eq:nigot}\ are % 式 (5.19)
\setcounter{equation}{18}
\begin{equation}\label{eq:pokodo}
\begin{aligned}\partial_{\bar y}(A_{\bar
z}+t'\phi_{\bar z})&=0\\ \partial_y(A_{\bar
z}-(t'')^{-1}\phi_{\bar z})&=0 .\\\end{aligned}
\end{equation}
 This determines the two
complex structures: % 式 (5.20)
\setcounter{equation}{19}
\begin{equation}\label{eq:plombob}
(w_+,w_-)=(-t',(t'')^{-1}).
\end{equation}

Now we can determine how $S$-duality acts on the full family
$\Bbb{CP}^1_h\times \Bbb{CP}^1_h$.  We learned in eqn. \eqref{eq:duco}\ that
a duality transformation $\left(\begin{matrix}a & b \\ c & d\end{matrix}\right)\in
SL(2,\Bbb{Z})$ (or its analog for a gauge group that is not simply-laced)
acts on $t$ by % 式 (5.21)
\setcounter{equation}{20}
\begin{equation}\label{eq:utlo}
t\to t{c\tau+d\over
|c\tau+d|}.
\end{equation}
  The same reasoning shows that $t'$ and $t''$
transform in precisely the same way.  {}From this and \eqref{eq:plombob}, it
follows that the action of $S$-duality on $\Bbb{CP}^1_h\times
\Bbb{CP}^1_h$ is % [WARN] 公式 \eqn\eqref{eq:yfog} 无编号或重复
\begin{equation*}
\tag*{(5.22)}
\begin{aligned}w_+ & \to w_+{c \tau+d\over
|c\tau+d|}\\
                   w_-  & \to w_-{|c\tau+d|\over c
                   \tau+d}.\\\end{aligned}
\end{equation*}

An important special case is that the $B$-model in complex
structure $I$, which corresponds to $(w_+,w_-)=(0,0)$, is
completely invariant under duality transformations. The $A$-model
in complex structure $I$, which corresponds to
$(w_+,w_-)=(0,\infty)$, is likewise invariant. This and other
statements about two-dimensional TQFT's that are not on the
distinguished family $\Bbb{CP}^1_g$ will be useful auxiliary
tools, but as will become clear, the geometric Langlands program
is really a statement or collection of statements about the
distinguished family.

\par\medskip{\it Dependence On The Metric}

In general, a two-dimensional TQFT with target $\MH(G,C)$ will
depend on the complex structure of $C$, because this influences
the hyper-Kahler structure of $\MH(G,C)$.  However, if a
two-dimensional TQFT, such as those parametrized by
$\Bbb{CP}^1_g$,  descends from a four-dimensional one, it must be
independent of the complex structure on $C$.  After all, the
four-dimensional TQFT on $M=\Sigma\times C$ is independent of the
metric on $M$ and in particular on $C$.

In some cases, we can see directly that these models are
independent of the complex structure of $C$.   At $t=\pm i$, we
get the $B$-model in complex structure $\pm J$; this complex
structure is independent of the metric of $C$, as we observed in
section 第~\ref{sec:mhhyper}~节. For real $t$, we get the $A$-model in a complex
structure that is a linear combination of complex structures $I$
and $K$. The $A$-model is determined by the corresponding
symplectic structure, which is a linear combination of symplectic
structures $\omega_I$ and $\omega_K$; we  found these in section
第~\ref{sec:mhhyper}~节\ to be independent of the metric of $C$.

To go farther, we will use generalized complex geometry.
